\documentclass[10pt,a4paper]{amsart}
\usepackage[margin=19mm]{geometry}
\usepackage{amsmath,amssymb,mathtools}
\usepackage[scr=rsfs,cal=euler]{mathalfa}
\usepackage{xfrac}
\usepackage[unicode,hidelinks,bookmarks=false]{hyperref}
\newcommand{\ie}{i.e.\ }
\newcommand{\cf}{cf.\ }
\newcommand{\resp}{respectively\ }
\newcommand{\Subsection}{\subsection}
\providecommand{\nobreakdash}{\nobreak\hbox{-}}
\setlength{\emergencystretch}{2em}
\multlinegap0pt
\usepackage{amssymb}
\usepackage{mathtools}
\newtheorem{lemma}{Lemma}
\newtheorem{proposition}{Proposition}
\newcommand{\Besov}{B^0_{\infty,1}}
\newcommand{\Log}{\operatorname{Log}}
\newcommand{\T}{\mathbb T}
\newcommand{\calE}{\mathcal E}
\newcommand{\calR}{\mathcal R}
\newcommand{\divv}{\operatorname{div}}
\newcommand{\eps}{\varepsilon}
\newcommand{\norm}[1]{\left\lVert #1\right\rVert}
\title{Pressure Quotients and Endpoint Velocity-Clock Criteria for Non-Diffusive Viscoelastic Flows}
\author{Sai Peng}
\date{Source: arXiv:2606.25258v4 (CC BY 4.0)}
\begin{document}
\maketitle

\noindent\emph{Source and licence.} Pressure Quotients and Endpoint Velocity-Clock Criteria for Non-Diffusive Viscoelastic Flows. Version arXiv:2606.25258v4 (CC BY 4.0), distributed by arXiv under the Creative Commons Attribution 4.0 International licence, http://creativecommons.org/licenses/by/4.0/, as recorded in the arXiv licence metadata for this exact version and shown on the abstract page https://arxiv.org/abs/2606.25258v4. Full licence notice: \texttt{licenses/arxiv-2606.25258-CC-BY-4.0.txt}.\par\smallskip

\noindent\textit{Editorial scope.} Counted unit: the article's proposition prop:physical-L2-estimate (source0.tex lines 2105--2128). The article's complete original proof of it is delivered with it (source0.tex lines 2130--2267, one \texttt{proof} environment of 138 lines). No mathematics is written, repaired or re-derived: the statement and the proof are the article's own lines, in the article's own notation, and no retained line is substituted or re-flowed.  Retained uncounted as the source-local context the proof needs: lines 1931--1939; lines 2015--2068.  Nothing was omitted from any retained span, so the package carries no omission record.  No retained line contains a citation, so the wrapper carries no bibliography and no bibliography entry.  No retained line contains a figure environment, an image-inclusion command or drawing code, so no image is delivered and the package declares no asset dependency.  Editorial formatting only, in addition: an AMS \texttt{amsart} preamble that restates the article's own declarations and macros used by the retained lines, namely the environments \texttt{lemma}, \texttt{proposition} on separate counters, together with the standard packages those lines need.  The retained environments print in this document as Lemma 1, Proposition 1 in order of appearance, whereas the article prints the same items with the numbering of its own full document.  The complete native source of the article as distributed in the arXiv e-print for this exact version is bundled unchanged as \texttt{sources/203588/source0.tex}.
\begin{align}
  \partial_tu+u\cdot\nabla u-\nu\Delta u+\nabla\pi
    &=\alpha\divv Y, \label{eq:physical-u}\\
  (\partial_t+u\cdot\nabla)a+\lambda^{-1}(a-1)&=S(u):Y,
    \label{eq:physical-a}\\
  (\partial_t+u\cdot\nabla)Y+\lambda^{-1}Y
    &=2aS(u)+[\nabla u\,Y+Y(\nabla u)^T]^\circ .
    \label{eq:physical-Y}
\end{align}

\begin{lemma}[Endpoint transport and physical commutator bounds]
\label{lem:physical-commutator}
Let \(m\ge3\) and \(0<\eps<1\).  For smooth fields on \(\T^2\) and
\(\divv u=0\),
\[
  \left|
  \sum_{|\beta|\le m}
  \int \partial^\beta f:
  [\partial^\beta,u\cdot\nabla]f\,dx
  \right|
  \le C\norm{\nabla u}_{\Besov}\norm{f}_{H^m}^2.
  \label{eq:pointwise-transport-commutator}
\]
Moreover, if \(w\in L^\infty\), then
\[
\begin{aligned}
  \left|
  \sum_{|\beta|\le m}
  \int w\,\partial^\beta f:
  [\partial^\beta,u\cdot\nabla]f\,dx
  \right|
  &\le C\norm{w}_{L^\infty}
  \norm{\nabla u}_{\Besov}\norm{f}_{H^m}^2 \\
  &\quad
  +C\norm{w}_{L^\infty}\norm{f}_{H^{1+\eps}}
    \norm{u}_{H^{m+1}}\norm{f}_{H^m}.
\end{aligned}
  \label{eq:weighted-transport-commutator}
\]
Moreover, for \(|\beta|\le m\),
\[
  \norm{[\partial^\beta,a]\nabla u}_{L^2}
  \le C\norm{a}_{H^{1+\eps}}\norm{u}_{H^{m+1}}
       +C\norm{\nabla u}_{\Besov}\norm{a}_{H^m},
  \label{eq:kp-a-u}
\]
and the same estimate holds with \(a\) replaced by any component of \(Y\).  In
particular,
\[
\begin{aligned}
  &\norm{[\partial^\beta,Y]\nabla u}_{L^2}
  +\norm{[\partial^\beta,\nabla u]Y}_{L^2}  \\
  &\qquad\le
  C\norm{Y}_{H^{1+\eps}}\norm{u}_{H^{m+1}}
  +C\norm{\nabla u}_{\Besov}\norm{Y}_{H^m} .
\end{aligned}
  \label{eq:kp-Y-u}
\]
The same right-hand side controls
\[
  \norm{\partial^\beta\bigl((\nabla u)Y+Y(\nabla u)^T\bigr)}_{L^2}
\]
after subtracting any explicitly displayed top-order term.
\end{lemma}

\begin{proposition}[Oldroyd--B active-deviatoric endpoint estimate]
\label{prop:physical-L2-estimate}
Let \(m\ge3\), \(0<\eps<1\), and suppose that
\[
  0<c_0I\le A(t,x)\le C_0I
\]
on \([0,T]\).  Define
\[
  \calE_m(t)=\norm{u}_{H^m}^2+\norm{a-1}_{H^m}^2
  +\frac{\alpha}{4}\sum_{|\beta|\le m}\int a^{-1}|\partial^\beta Y|^2\,dx .
\]
Then
\[
\begin{aligned}
  \frac{d}{dt}\calE_m
  &+c\nu\norm{u}_{H^{m+1}}^2
   +c\lambda^{-1}\bigl(\norm{a-1}_{H^m}^2+\alpha\norm{Y}_{H^m}^2\bigr) \\
  &\le
  C_{c_0,C_0}\left(1+\norm{\nabla u}_{\Besov}
  +\norm{\Log A}_{H^{1+\eps}}^2\right)\calE_m .
\end{aligned}
\label{eq:physical-L2-energy}
\]
\end{proposition}

\begin{proof}
Write \(D_t=\partial_t+u\cdot\nabla\).  The compact spectral window gives
uniform upper and lower bounds for \(a\) and \(a^{-1}\), so \(\calE_m\) is
equivalent to
\(\norm{u}_{H^m}^2+\norm{a-1}_{H^m}^2+\alpha\norm{Y}_{H^m}^2\).

Apply \(\partial^\beta\), \(|\beta|\le m\), to
\eqref{eq:physical-u} and test by \(\partial^\beta u\).  Since \(\divv u=0\),
the transport term is a commutator and the pressure term vanishes.  The elastic
term is
\[
  \alpha\int \partial^\beta u\cdot\partial^\beta\divv Y\,dx
  =-\alpha\int \partial^\beta Y:S(\partial^\beta u)\,dx .
  \label{eq:principal-velocity-coupling}
\]
The remaining velocity terms are bounded by
\[
  C\norm{\nabla u}_{\Besov}\norm{u}_{H^m}^2
  -\nu\norm{\nabla\partial^\beta u}_{L^2}^2 .
\]

Next differentiate the anisotropic equation \eqref{eq:physical-Y}.  After
commuting \(\partial^\beta\) with the material derivative,
\[
\begin{aligned}
  D_t\partial^\beta Y+\lambda^{-1}\partial^\beta Y
  &=2aS(\partial^\beta u)
    +2[\partial^\beta,a]S(u)  \\
  &\quad+\partial^\beta\bigl([
      \nabla u\,Y+Y(\nabla u)^T]^\circ\bigr)
    -[\partial^\beta,u\cdot\nabla]Y .
\end{aligned}
\]
Test this identity by \((\alpha/2)a^{-1}\partial^\beta Y\).  Because
\(\divv u=0\),
\[
\begin{aligned}
  &\frac{\alpha}{4}\frac{d}{dt}\int a^{-1}|\partial^\beta Y|^2\,dx
  +\frac{\alpha}{2\lambda}\int a^{-1}|\partial^\beta Y|^2\,dx  \\
  &=\alpha\int \partial^\beta Y:S(\partial^\beta u)\,dx
    +\calR_\beta^{a}+\calR_\beta^{Y}
    +\calR_\beta^{tr}+\calR_\beta^{w},
\end{aligned}
\label{eq:weighted-Y-identity}
\]
where
\[
  \calR_\beta^{a}
  =\alpha\int a^{-1}\partial^\beta Y:
  [\partial^\beta,a]S(u)\,dx,
\]
\(\calR_\beta^Y\) contains the quadratic stretching term
\(\partial^\beta([\nabla u\,Y+Y(\nabla u)^T]^\circ)\),
\(\calR_\beta^{tr}\) contains the transport commutator, and
\[
  \calR_\beta^{w}
  =\frac{\alpha}{4}\int D_t(a^{-1})|\partial^\beta Y|^2\,dx
\]
is the contribution of the time-dependent weight.  The first term on the right-hand side of the weighted identity cancels exactly
with the elastic term from the velocity equation:
\[
  -\alpha\int \partial^\beta Y:S(\partial^\beta u)\,dx
  +\alpha\int \partial^\beta Y:S(\partial^\beta u)\,dx=0.
\]
This is the top-order cancellation identity.  It occurs after the isotropic
stress has been removed into the pressure and before the use of Young's
inequality.  Consequently the remaining estimates have to control only
commutators, the time derivative of the weight \(a^{-1}\), the scalar equation
for \(a\), and lower-order stretching terms.  This is the source of the
\(L^2_t\) logarithmic coefficient.

We now bound the remainders.  For the coefficient commutator,
Lemma~\ref{lem:physical-commutator} gives, for every \(\delta>0\),
\[
\begin{aligned}
  |\calR_\beta^{a}|
  &\le C\norm{\partial^\beta Y}_{L^2}
       \norm{[\partial^\beta,a]S(u)}_{L^2}  \\
  &\le \delta\nu\norm{u}_{H^{m+1}}^2
      +C_\delta\norm{a}_{H^{1+\eps}}^2\norm{Y}_{H^m}^2
      +C\norm{\nabla u}_{\Besov}\calE_m .
\end{aligned}
\label{eq:a-commutator-remainder}
\]
The same lemma gives
\[
  |\calR_\beta^{Y}|
  \le \delta\nu\norm{u}_{H^{m+1}}^2
      +C_\delta\norm{Y}_{H^{1+\eps}}^2\norm{Y}_{H^m}^2
      +C\norm{\nabla u}_{\Besov}\calE_m .
  \label{eq:Y-stretching-remainder}
\]
The transport commutator is controlled by the weighted transport estimate in
Lemma~\ref{lem:physical-commutator}.  The part with a top derivative on \(u\)
is absorbed by viscosity:
\[
  |\calR_\beta^{tr}|
  \le \delta\nu\norm{u}_{H^{m+1}}^2
      +C_\delta\norm{Y}_{H^{1+\eps}}^2\norm{Y}_{H^m}^2
      +C_{c_0,C_0}\norm{\nabla u}_{\Besov}\norm{Y}_{H^m}^2 .
  \label{eq:Y-transport-remainder}
\]
Finally, from \eqref{eq:physical-a},
\[
  D_t(a^{-1})=\lambda^{-1}a^{-2}(a-1)-a^{-2}S(u):Y .
\]
Hence, using the compact spectral window and
\(\norm{\nabla u}_{L^\infty}\le C\norm{\nabla u}_{\Besov}\),
\[
  |\calR_\beta^w|
  \le C_{c_0,C_0}\left(1+\norm{\nabla u}_{\Besov}\right)
      \norm{Y}_{H^m}^2 .
  \label{eq:weight-remainder}
\]

It remains to estimate the scalar equation.  Applying \(\partial^\beta\) to
\eqref{eq:physical-a}, testing by \(\partial^\beta(a-1)\), and using
Lemma~\ref{lem:physical-commutator} gives
\[
\begin{aligned}
  \frac12\frac{d}{dt}\norm{a-1}_{H^m}^2
  +\lambda^{-1}\norm{a-1}_{H^m}^2
  &\le \delta\nu\norm{u}_{H^{m+1}}^2  \\
  &\quad+C_\delta\norm{Y}_{H^{1+\eps}}^2\norm{a-1}_{H^m}^2
  +C\norm{\nabla u}_{\Besov}\calE_m .
\end{aligned}
\label{eq:a-high-order}
\]

On the compact spectral window, the smooth maps
\(B=\Log A\mapsto a\) and \(B\mapsto Y\) satisfy
\[
  \norm{a}_{H^{1+\eps}}+\norm{Y}_{H^{1+\eps}}
  \le C_{c_0,C_0}\bigl(1+\norm{\Log A}_{H^{1+\eps}}\bigr).
  \label{eq:window-log-composition}
\]
Choose \(\delta>0\) small enough that the velocity terms in the displayed commutator, stretching, and scalar estimates are absorbed by the viscous dissipation.  Summing over \(|\beta|\le m\), using the compact-window composition estimate, and using the weighted equivalence of \(\calE_m\) gives the claimed differential inequality.
\end{proof}

\end{document}
